% TOP_PROBLEM: {"id": 30007281, "title": "Constructible tame eigensheaves with regular monodromy", "category_id": 1, "category": {"id": 1, "name": "number_theory", "display_name": "Number Theory", "description": "Properties of integers, prime numbers, Diophantine equations.", "slug": "number-theory", "order_index": 1, "created_at": "2026-07-31T15:26:25.670Z"}, "status": "open"}
% FIRST_POSTED: 2026-10-06
% ENTRY_KIND: statement_only
% !TeX program = lualatex
\documentclass[11pt]{article}
\usepackage[margin=25mm]{geometry}
\usepackage{fontspec}
\setmainfont{Latin Modern Roman}
\usepackage{amsmath,amssymb,mathtools}
\usepackage{unicode-math}
\setmathfont{Latin Modern Math}
\usepackage{xurl,hyperref}
\hypersetup{colorlinks=true,urlcolor=blue}
\setcounter{secnumdepth}{0}
\setlength{\parindent}{0pt}
\setlength{\parskip}{5pt}
\setlength{\emergencystretch}{3em}
\newcommand{\R}{\mathbb R}
\newcommand{\N}{\mathbb N}
\newcommand{\Z}{\mathbb Z}
\newcommand{\Q}{\mathbb Q}
\newcommand{\C}{\mathbb C}
\newcommand{\D}{\mathbb D}
\newcommand{\F}{\mathbb F}
\newcommand{\sref}[2]{\hyperlink{source:#1:#2}{[S#2]}}
\newcommand{\cataloguescope}[1]{#1}
\begin{document}
\section*{Constructible tame eigensheaves with regular monodromy}\label{30007281}
\subsection{Definitions and mathematical statement}
\paragraph{Statement recorded in the supplied source.}
Let X be a smooth projective geometrically connected curve of genus at least two over an algebraic closure of a finite field of characteristic p, with a finite-field rational marked point x, and let G be split simple simply connected. Assume that the Lie algebra admits an invariant nondegenerate symmetric form nondegenerate on every Levi center, Chevalley restriction is an isomorphism for every Levi, semisimple Lie-element scheme centralizers are Levi subgroups, and every nilpotent Lie element over every extension belongs to the Lie algebra of the unipotent radical of an extension-defined parabolic. For every Zariski-dense geometric dual-group parameter on X minus x over a finite extension of $Q_\ell $, ell different from p, killing wild inertia and having tame monodromy $\exp (t_\ell $ N) for regular nilpotent N, is there a nonzero locally constructible perverse ell-adic sheaf on G-bundles with Borel reduction at x, with parabolic nilpotent singular support and coherent all-leg Hecke eigenisomorphisms, including fusion at collisions? No arithmetic descent of the parameter is assumed.
\subsection{Short English statement}
Let X be a smooth projective geometrically connected curve of genus at least two over an algebraic closure of a finite field of characteristic p, with a finite-field rational marked point x, and let G be split simple simply connected. Assume that the Lie algebra admits an invariant nondegenerate symmetric form nondegenerate on every Levi center, Chevalley restriction is an isomorphism for every Levi, semisimple Lie-element scheme centralizers are Levi subgroups, and every nilpotent Lie element over every extension belongs to the Lie algebra of the unipotent radical of an extension-defined parabolic. For every Zariski-dense geometric dual-group parameter on X minus x over a finite extension of $Q_\ell $, ell different from p, killing wild inertia and having tame monodromy $\exp (t_\ell $ N) for regular nilpotent N, is there a nonzero locally constructible perverse ell-adic sheaf on G-bundles with Borel reduction at x, with parabolic nilpotent singular support and coherent all-leg Hecke eigenisomorphisms, including fusion at collisions? No arithmetic descent of the parameter is assumed.
\subsection{Sources}
\begin{itemize}
\item[S1] Original question source or mathematical context.
\url{https://doi.org/10.4171/owr/2016/20}.
\item[S2] Original question source or mathematical context.
\url{https://doi.org/10.4171/owr/2021/39}.
\item[S3] OpenAI, manuscript citation page. Source of the formulation motivating this local problem.
\url{https://github.com/openai/math/blob/adc7f1241b42e322a6451854ab7e4b4c146bf78a/preprints/Constructible-tame-Hecke-eigensheaves-in-positive-characteristic-October-5-2026/README.md}.
\end{itemize}
\end{document}
